% ---
% title: Allocating capital
% tags: finance, stars-and-bars
% ---

\begin{problem}
	\begin{context}
		\emph{Capital allocation} means deciding how much money goes into
		each investment. Real opportunities often have a minimum ticket size (a
		fund or a project accepts only investments above a threshold).
	\end{context}
	We have 20 thousand euros to invest in 4 opportunities, in whole units of
	1000. The minimal investments are 2, 2, 3 and 4 thousand. How many
	investment strategies are there if (a) we invest in each opportunity;
	(b) we invest in at least 3 of the 4? (All 20 units are invested.)
\end{problem}

\begin{slides}
	\small
	\exstatement
	\pause
	\sol (a) First pay the minimums: $2+2+3+4=11$ units. The remaining $9$
	identical units go to 4 boxes freely (stars and bars, $9$ stars, $3$ bars):
	$\binom{9+3}{3}=\binom{12}{3} = 220$.\\[1mm]
	(b) Add the cases where exactly one opportunity gets nothing:\par\smallskip
	{\centering
		\begin{tabular}{lccc}\toprule
			left out & minimums paid & left to divide & count               \\\midrule
			1st (2)  & $2+3+4=9$     & 11 among 3     & $\binom{13}{2}=78$  \\
			2nd (2)  & $2+3+4=9$     & 11 among 3     & $78$                \\
			3rd (3)  & $2+2+4=8$     & 12 among 3     & $\binom{14}{2}=91$  \\
			4th (4)  & $2+2+3=7$     & 13 among 3     & $\binom{15}{2}=105$ \\\bottomrule
		\end{tabular}\par}\smallskip
	Total: $220+78+78+91+105 = 572$.
\end{slides}

\begin{solution}
	\idea The units of 1000 euros are identical, the opportunities are
	different: stars and bars. Minimums are handled by paying them first.

	\step{(a)} Pay the minimums: $2+2+3+4=11$ units. The remaining $20-11=9$
	units can be divided freely among the 4 opportunities (some may get nothing
	extra): 9 stars and 3 bars,
	\[ \binom{9+3}{3}=\binom{12}{3}=220. \]

	\step{(b)} ``At least 3'' = ``all 4'' (part (a)) or ``exactly 3''. These
	cases are disjoint, so we add. If exactly one opportunity is left out, we pay
	the minimums of the other three and divide the rest among those three
	(2 bars):
	\begin{center}
		\begin{tabular}{lccc}\toprule
			left out    & minimums paid & left to divide & count               \\\midrule
			1st (min 2) & $2+3+4=9$     & 11 among 3     & $\binom{13}{2}=78$  \\
			2nd (min 2) & $2+3+4=9$     & 11 among 3     & $\binom{13}{2}=78$  \\
			3rd (min 3) & $2+2+4=8$     & 12 among 3     & $\binom{14}{2}=91$  \\
			4th (min 4) & $2+2+3=7$     & 13 among 3     & $\binom{15}{2}=105$ \\\bottomrule
		\end{tabular}
	\end{center}
	Total:
	\[ 220+78+78+91+105=572. \]
	\begin{answer}
		(a) $220$ \qquad (b) $572$
	\end{answer}
	\begin{remark}
		\textbf{Common mistake:} in the ``left out'' cases still paying the left-out
		opportunity's minimum, or dividing among 4 boxes instead of 3. A left-out
		opportunity gets exactly $0$.
	\end{remark}
\end{solution}
